Set $g(x)=f(x)-x^2$. The equation becomes $g(x+y)=g(x)+g(y)$, so the complete family is $\boxed{f(x)=x^2+g(x)}$ for arbitrary additive $g$. Conversely every such function works by expanding $(x+y)^2$. In particular $f=x^2+cx$ works for any real constant $c$, so uniqueness already fails among polynomials. Without regularity assumptions, additive functions need not be linear. If continuity is imposed, additivity gives $g(q)=qg(1)$ on rational $q$ and then by limits on every real number.
